\subsection{代表域}

定义 1．——设 A 是一个局部环。A 的代表域是 A 的任一子域 K，使 A 到 \(\kappa_{\mathrm{A}}\) 的典范同态诱导出从 K 到 \(\kappa_{\mathrm{A}}\) 的同构（换言之，满足 \(\mathrm{A} = \mathrm{K} + \mathfrak{m}_{\mathrm{A}}\)）。

A 的代表域只有在 A 有特征时才可能存在。当 A 分离且完备时，这个条件也是充分的。更确切地说，有如下定理：

定理 1．——设 A 是特征为 p 的分离且完备局部环。

\begin{enumerate}
\def\labelenumi{\alph{enumi})}
\item
  假设 \(p = 0\)，并令 \((x_{\lambda})_{\lambda \in \Lambda}\) 为 A 中的一个元素族，其模 \(\mathfrak{m}_{\mathbf{A}}\) 的类构成 \(\kappa_{\mathbf{A}}\) 在 \(\mathbf{Q}\) 上的超越基。存在唯一的 A 的代表域包含元素 \(x_{\lambda}\)。
\item
  假设 \(p \neq 0\)。令 \((x_{\lambda})_{\lambda \in \Lambda}\) 为 A 中的一个元素族，其模 \(\mathfrak{m}_{\mathbf{A}}\) 的类构成 \(\kappa_{\mathbf{A}}\) 的 p-基（A，V，p. 95）。存在唯一的 A 的代表域包含元素 \(x_{\lambda}\)。它是 A 的一个科恩子环。
\end{enumerate}

设 \(p = 0\)，于是 A 是 Q-代数。由于 \(\kappa_{\mathrm{A}}\) 在 Q 上的每个超越基都是可分的，断言 \(a)\) 由第 1 号的命题 1 应用于 \(k_0 = \mathbf{Q}\)、\(\mathbf{K} = \kappa_{\mathrm{A}}\) 的情形得出。

现在设 \(p \neq 0\)。于是有 \(p1_{\mathbf{A}} = 0\)，且 A 的每个科恩子环 C 都满足 \(p\mathbf{C} = 0\)。换言之，A 的代表域与科恩子环这两个概念相同。断言 \(b)\) 随即由 § 2 第 2 号的定理 1 得出。

推论 1．——设 A 是分离且完备的局部环，其剩余域是 Q 的代数扩张。于是存在唯一的 A 的代表域。事实上，环 A 的特征为 0（第 1 号）。

推论 2．——设 A 是特征为 \(p \neq 0\) 的分离完备局部环。假设剩余域 \(\kappa_{\mathrm{A}}\) 是完美的。则存在唯一的 A 的代表域，即乘法代表元的集合。

推论 2 立即由定理 1 以及 § 2 第 4 号的命题 7 得出。

定理 2．——设 A 是含有一个域的维数为 d 的完备诺特局部环。设 K 是 A 的一个代表域，并令 m 为向量空间 \(\mathfrak{m}_{\mathrm{A}} / \mathfrak{m}_{\mathrm{A}}^{2}\) 在域 K 上的维数。

\begin{enumerate}
\def\labelenumi{\alph{enumi})}
\item
  存在理想 \(\mathfrak{a}\)，它属于环 \(\mathbf{K}[[\mathrm{T}_1,\dots,\mathrm{T}_m]]\)，使 K-代数 A 同构于 \(\mathbf{K}[[\mathrm{T}_1,\dots,\mathrm{T}_m]] / \mathfrak{a}\)。
\item
  存在 A 的一个子 K-代数 A'，同构于 \(K[[T_{1}, \ldots, T_{d}]]\)，且 A 是 A' 上的有限代数。
\item
  假设诺特局部环 A 是正则的，即 \(d = m\)。则存在从 A 到 \(\mathbf{K}[[\mathbf{T}_1, ..., \mathbf{T}_d]]\) 的 K-同构。
\end{enumerate}

令元素 \(t_{1},\ldots,t_{m}\) 属于 \(\mathfrak{m}_{\mathrm{A}}\)，其模 \(\mathfrak{m}_{\mathrm{A}}^{2}\) 的类生成 K-向量空间 \(\mathfrak{m}_{\mathrm{A}}/\mathfrak{m}_{\mathrm{A}}^{2}\)。根据 § 2 第 5 号的引理 3，存在从 \(\mathrm{K}[[\mathrm{T}_{1},\ldots,\mathrm{T}_{m}]]\) 到 A 的满射 K-同态，将 \(\mathrm{T}_{i}\) 映到 \(t_{i}\)，对 \(1\leqslant i\leqslant m\) 均如此。这证明了 \(a)\)。

同理，断言 \(b\) 由同上引理 4 以及 A 的极大正则序列的存在性得出（VIII，§ 3，第 2 号，定理 1）。

最后，断言 c) 正是 VIII，§ 5，第 2 号定理 1 的推论 3。

